\documentclass{article}
\usepackage{mathtools} 
\usepackage{fontspec}
\usepackage[UTF8]{ctex}
\usepackage{amsthm}
\usepackage{mdframed}
\usepackage{xcolor}
\usepackage{amssymb}
\usepackage{amsmath}


% 定义新的带灰色背景的说明环境 zremark
\newmdtheoremenv[
  backgroundcolor=gray!10,
  % 边框与背景一致，边框线会消失
  linecolor=gray!10
]{zremark}{说明}

% 通用矩阵命令: \flexmatrix{矩阵名}{元素符号}{行数}{列数}
\newcommand{\flexmatrix}[4]{
  \[
  #1 = \begin{pmatrix}
    #2_{11}     & #2_{12}     & \cdots & #2_{1#4}   \\
    #2_{21}     & #2_{22}     & \cdots & #2_{2#4}   \\
    \vdots      & \vdots      & \ddots & \vdots     \\
    #2_{#31}    & #2_{#32}    & \cdots & #2_{#3#4}
  \end{pmatrix}
  \]
}

% 简化版命令（默认矩阵名为A，元素符号为a）: \quickmatrix{行数}{列数}
\newcommand{\quickmatrix}[2]{\flexmatrix{A}{a}{#1}{#2}}

\begin{document}
\title{注释 9.4}
\author{张志聪}
\maketitle

\begin{zremark}
  例2引出处的推论（积分不等式的推广）：\\
  若$f$与$g$为$[a,b]$上的两个可积函数，且$f(x) \leq g(x), x \in [a,b]$，
  只要存在$x_0 \in [a,b]$，$f,g$在$x_0$处连续，且$f(x_0) < g(x_0)$，则有
  \begin{align*}
    \int_a^b f(x) dx < \int_a^b g(x) dx
  \end{align*}
\end{zremark}

\textbf{证明：}

提示：令$h(x) = g(x) - f(x)$，考虑函数$h(x)$的积分。

\begin{zremark}
  $f$在$[a, b]$上可积，若$f(x) \geq 0, x \in [a, b]$，则
  \begin{align*}
    F(x) = \int_a^x f(t) dt
  \end{align*}
  在$[a, b]$上单增。
\end{zremark}
以几何的角度来看，这个结论是显然的。

\textbf{证明：}

任意$x_1, x_2 \in [a, b]$（设$x_1 < x_2$），于是
\begin{align*}
  F(x_2) - F(x_1)
   & = \int_a^{x_2} f(t) dt - \int_a^{x_1} f(t) dt                            \\
   & = \int_a^{x_1} f(t) dt + \int_{x_1}^{x_2} f(t) dt - \int_a^{x_1} f(t) dt \\
   & = \int_{x_1}^{x_2} f(t) dt                                               
\end{align*}
由积分不等式可知，$\int_{x_1}^{x_2} f(t) dt \geq 0$，即
\begin{align*}
  F(x_2) - F(x_1) \geq 0
\end{align*}



\begin{zremark}
  复合函数的积分
\end{zremark}

\textbf{证明：}

todo

\end{document}